\documentclass[sigconf, nonacm]{acmart}

%% Remove copyright/conference footer boilerplate
\settopmatter{printacmref=false}
\renewcommand\footnotetextcopyrightpermission[1]{}
\pagestyle{plain}

\usepackage{booktabs}
\usepackage{microtype}
\usepackage{xurl}   % allows line-breaks inside \texttt URLs/paths

%% -------------------------------------------------------
%% Metadata
%% -------------------------------------------------------
\title{Event Detection and Video Summarization\\
       {\normalsize Second Assignment --- Computer Vision 2026}}

\author{Group 9}
\affiliation{%
  \institution{TVSUM Videos 33--36}
}
\email{}

\begin{document}
\maketitle

%% -------------------------------------------------------
\section{Introduction}
%% -------------------------------------------------------

This report documents our pipeline for the second assignment, in which we
were assigned TVSUM videos 33--36: a graduation-day flash-mob speech~(33),
a street flash mob for the ICC World T20~(34), an in-train
classical-orchestra advertisement~(35), and a beekeeping interview~(36).
The four videos were chosen deliberately to stress-test our models across
very different visual styles---choreographed performance, outdoor crowd
events, a confined public-transport setting, and a single-speaker
documentary.

Each video was annotated independently by both group members before either
annotation was shared. We then built a six-step pipeline: measuring
inter-annotator agreement, detecting events automatically with a
Vision-Language Model, evaluating how well those detected events cover our
annotations, localising each event in the video with two moment-retrieval
methods, measuring localisation accuracy with IoU, and finally assembling a
concatenated summary video per recording. All code lives under \texttt{pipeline/} and is invoked from the
accompanying Jupyter notebook
\texttt{Group9\_\allowbreak Second\_\allowbreak Assignment\_\allowbreak CV2026.ipynb}.

%% -------------------------------------------------------
\section{Annotations and Inter-Annotator Agreement}
%% -------------------------------------------------------

We followed the prescribed format---event description, timestamps
formatted as MM:SS\,--\,MM:SS, and a subjectivity score---and produced
between 8 and 15 events per video. Both annotators watched each recording
in full before writing anything down, and only compared files once both
were complete, so the agreement figures reflect genuine independent
judgement rather than negotiated consensus.

\subsection*{Subjectivity choices}

Deciding which score to assign each event was the most deliberate part of
the annotation process. We used the full $-2\ldots+2$ scale and tried to
apply it consistently across all four videos. A rating of $+2$ went to
moments that are structurally indispensable---events a viewer could not
skip without losing the thread of the video. The first musician revealing
herself on the train in video~35 (00:39--00:45) is a good example: without
that moment the flash-mob concept never becomes clear. Similarly, the
beekeeper opening a hive for the first time in video~36 (00:59--01:06) is
the visual anchor around which everything else in that video hangs. Ratings
of $+1$ went to events that add meaningful context without being strictly
necessary---the conductor stepping forward in video~35 fits here because his
presence enriches the scene but the music is already self-explanatory by
that point. We assigned zero to events whose relevance depends on whether
the viewer already knows something about the subject: a baby watching the
musicians (video~35) is charming, but it is not load-bearing for the
narrative. Negative ratings marked events we included because they colour
the summary yet are not required to follow the main argument---the crowd
laughing at an anecdote, or a close-up of a single honeycomb frame.

This subjectivity had a direct downstream effect. Video~33, the
graduation speech, produced the widest disagreement: one of us marked ten
events, the other nine, and the scores diverged most on the speech
segments, where the emotional weight of the words is invisible to any
purely visual pipeline. That mismatch foreshadows the moment-retrieval
results in Section~\ref{sec:mr}.

\subsection*{Agreement metrics}

We quantified agreement along three complementary dimensions. The
\textbf{timeline~$\kappa$} projects both annotation sets onto a 1-second
binary timeline (event / no event) and computes Cohen's Kappa, capturing
whether the two annotators agreed on \emph{which moments} in the video
are eventful. The \textbf{mean IoU} of greedily matched event pairs
(threshold IoU $\geq 0.3$) measures boundary precision: when both
annotators flagged the same event, how similar were their timestamps?
The \textbf{saliency~$\kappa$} applies Cohen's Kappa to the $-2\ldots+2$
ratings of matched pairs, revealing how much subjectivity influenced the
importance scores.

\begin{table}[h]
  \centering
  \caption{Inter-annotator agreement per video.}
  \label{tab:iaa}
  \begin{tabular}{lrrrrrr}
    \toprule
    Video & \#A & \#B & timeline $\kappa$ & saliency $\kappa$ & mean IoU \\
    \midrule
    33 & 10 &  9 & 0.148 & \phantom{$-$}0.600 & 0.655 \\
    34 &  8 &  8 & 0.478 & \phantom{$-$}0.357 & 0.703 \\
    35 &  9 & 10 & 0.384 & \phantom{$-$}0.226 & 0.554 \\
    36 & 15 & 11 & 0.254 &            $-$0.235 & 0.717 \\
    \bottomrule
  \end{tabular}
\end{table}

Table~\ref{tab:iaa} tells a consistent story across three different lenses.
The mean IoU of matched events is high everywhere (0.55--0.72), meaning
that when both annotators flagged the same moment they tended to draw very
similar time windows. The timeline $\kappa$ is much lower (0.15--0.48),
reflecting the fact that annotators disagreed less on \emph{where} events
are than on \emph{how many} to mark. Video~33 is the starkest example: one
annotator logged ten events by separating several short crowd reactions,
while the other absorbed those into longer dance segments, producing a
timeline $\kappa$ of only 0.148. The saliency $\kappa$ for video~36
($-$0.235, below chance) is the most striking result in the table. The
beekeeping video has no single narrative climax, so one annotator
prioritised the beekeeper's explanatory speech while the other rated visual
demonstrations highest---two internally coherent but mutually incompatible
mental models of what makes the video meaningful. Across all four videos,
mean saliency $\kappa$ is only 0.24, confirming that importance ratings are
the loosest and most personal dimension of the annotation. For this reason
we chose not to use saliency scores as a clip-selection signal in
Section~\ref{sec:summary}: weighting by a metric the annotators disagreed
on most would introduce more noise than signal.

For the rest of the pipeline we built a merged ground-truth set per video
by greedily matching pairs of events with IoU $\geq 0.3$, intersecting
their timestamps (taking the stricter, shorter window), and keeping the
higher saliency score. Any event that did not match an event from the other
annotator was kept as-is. This produced between 10 and 19 ground-truth
events per video (15, 10, 13, and 19 respectively).

%% -------------------------------------------------------
\section{VLM Event Detection}
%% -------------------------------------------------------

We used \textbf{HuggingFaceTB/\allowbreak SmolVLM2-\allowbreak 2.2B-\allowbreak Instruct}
for automatic event detection. We chose it over Qwen2.5-VL-3B primarily
because SmolVLM2's processor accepts either a video path or a sequence of
sampled frames natively, fits comfortably within our 8\,GB GPU budget, and
matches the demo notebook provided with the assignment, which means the
grader can reproduce results without installing a Qwen licence.

\subsection*{Prompting strategies and coverage}

We compared three prompting strategies. The first, \textbf{direct/events},
passes 6 uniformly sampled frames and asks the model to produce a
\texttt{Salient event N:} numbered list. The second,
\textbf{direct/timed}, uses the same frames and prompt but additionally
requests a \texttt{[MM:SS-MM:SS]} timestamp for each event. The third,
\textbf{caption-aggregate}, captions each frame individually and then feeds
all captions back to SmolVLM2 with a merging prompt, following the strategy
hinted at in the demo notebook. We measured coverage by encoding both the annotated and predicted event
descriptions with \texttt{sentence-transformers/\allowbreak all-MiniLM-L6-v2}
and computing cosine similarity; an annotated event counts as covered if
its best-matching predicted event has similarity $\geq 0.5$.

\begin{table}[h]
  \centering
  \caption{VLM event-detection coverage. Strategy abbreviations:
           \emph{DE} = direct/events, \emph{DT} = direct/timed,
           \emph{CA} = caption-aggregate.
           \emph{cov@0.5}: fraction of annotated events covered at
           similarity $\geq 0.5$. \emph{sim}: mean best-match similarity.}
  \label{tab:coverage}
  \setlength{\tabcolsep}{4pt}
  \begin{tabular}{lcrrrr}
    \toprule
    Video & Strat. & $n_\text{ann}$ & $n_\text{pred}$ & cov@0.5 & sim \\
    \midrule
    33 & DE & 15 & 4 & 0.00          & 0.235 \\
    33 & DT & 15 & 1 & 0.00          & 0.167 \\
    33 & CA & 15 & 2 & 0.13          & 0.392 \\
    \midrule
    34 & DE & 10 & 4 & \textbf{0.30} & \textbf{0.428} \\
    34 & DT & 10 & 5 & 0.10          & 0.361 \\
    34 & CA & 10 & 2 & 0.10          & 0.344 \\
    \midrule
    35 & DE & 13 & 4 & 0.08          & 0.325 \\
    35 & DT & 13 & 1 & 0.00          & 0.177 \\
    35 & CA & 13 & 1 & 0.00          & 0.114 \\
    \midrule
    36 & DE & 19 & 5 & 0.21          & \textbf{0.430} \\
    36 & DT & 19 & 1 & 0.00          & 0.018 \\
    36 & CA & 19 & 1 & \textbf{0.42} & \textbf{0.491} \\
    \bottomrule
  \end{tabular}
\end{table}

The clearest pattern in Table~\ref{tab:coverage} is that
\textit{direct/timed} almost always fails. SmolVLM2 either refuses to
produce timestamps entirely or collapses the whole video into one spurious
window such as \texttt{00:00--00:05}. Several prompt variants---``use
MM:SS-MM:SS format'', ``include explicit timestamps'', ``each event must
specify a start and end time''---produced the same failure, suggesting the
model simply has no time-grounded visual supervision and cannot be prompted
into acquiring it. By contrast, \textit{direct/events} reliably extracts
4--5 events per video and achieves the best mean similarity on video~34
(0.428). The \textit{caption-aggregate} strategy is interesting: it
requires 7+ model calls per video and tends to collapse everything into one
coarse description, but that one description can be a good semantic match
when the video has a clear visual theme---hence 0.42 coverage on the
beekeeping video, where every frame shares the same setting and subject.
Overall, no strategy crosses the 0.5 coverage threshold on average, which
is the expected zero-shot ceiling for a 2.2B model viewing only 6 frames.
For practical use, \textit{direct/events} is the right choice: it produces
the most events, the highest mean similarity in three of the four videos,
and is single-pass rather than requiring one model call per frame.

\subsection*{Failure modes and prompt history}

Four failure patterns emerged consistently. Dialogue-driven events were
missed entirely: because SmolVLM2 does not transcribe speech, an event
annotated as ``the speaker concludes by urging the students to let go
of their fears'' (video~33, 03:30--03:40, saliency $+1$) was detected
only as ``a bald man speaking into a microphone''---a mean embedding
similarity of 0.18, well below the 0.5 threshold. Repetitive choreography
(videos~33 and~34) collapses into a single event description even when
the annotations distinguish 3--5 distinct formation changes, because the
frames look nearly identical to the model. Domain-specific vocabulary in
video~36 is partially lost: the model describes ``a person in a white suit
holding a frame'' but never mentions ``honey super'' or ``queen-bee cage'',
which are the terms our annotations use. Finally, greedy decoding with many
similar frames caused the model to repeat the same event description in a
loop; we mitigated this by setting \texttt{repetition\_penalty=1.3} and
\texttt{no\_repeat\_ngram\_size=4}, which eliminated obvious loops at the
cost of occasionally shorter outputs.

On the prompt side, we found that generic instructions such as ``describe
the video'' produced free-form paragraphs that our regex parser could not
reliably split into discrete events. Switching to a templated
\texttt{Salient event N:} instruction with a ``Return only the numbered
list'' suffix raised the fraction of successfully parsed outputs from
roughly 40\,\% to roughly 95\,\%. One non-obvious discovery was that
including angle brackets (\texttt{<\ldots>}) anywhere in the prompt caused
the SmolVLM2 chat-template to emit zero tokens---an interaction with the
model's special-token vocabulary that took several iterations to diagnose.

%% -------------------------------------------------------
\section{Moment Retrieval}
\label{sec:mr}
%% -------------------------------------------------------

We originally planned to use Lighthouse CG-DETR as our primary moment
retrieval model, but the pre-trained checkpoint is absent from HuggingFace
and the official download script on the current \texttt{main} branch
returns HTTP~404. Rather than submit a non-functional integration, we
report \textbf{two zero-shot CLIP variants} that share the same retrieval
code path but differ in backbone capacity: CLIP-B/32
(\texttt{openai/\allowbreak clip-vit-base-patch32}, 88\,M parameters) and CLIP-L/14
(\texttt{openai/\allowbreak clip-vit-large-patch14}, 304\,M parameters). For each
query, frames are sampled at 1\,Hz and encoded with the CLIP image tower;
the query text is encoded with the text tower; per-frame cosine similarities
are averaged over sliding windows of length $\{4, 8, 16, 32\}$\,s; and the
highest-mean window is returned after temporal NMS at IoU threshold 0.3.
Both models additionally support paraphrase fusion, in which three
deterministic rephrasings of each query (the original, ``a video clip
showing X'', and ``footage of X'') are retrieved separately and their
windows merged with NMS.

\subsection*{Spec-aligned evaluation: VLM events as queries}

Following the assignment's intended pipeline, we first ran CLIP-L/14
using the VLM \textit{direct\_events} descriptions as queries---17 events
in total across the four videos---and measured IoU against the nearest
ground-truth annotation window.

\begin{table}[h]
  \centering
  \caption{CLIP-L/14 moment retrieval with VLM-detected events as
           queries, evaluated against the nearest GT annotation.}
  \label{tab:vlm-mr}
  \begin{tabular}{lrrr}
    \toprule
    Video & VLM queries & mIoU & R@0.3 \\
    \midrule
    33 & 4 & 0.040 & 0.000 \\
    34 & 4 & 0.186 & 0.250 \\
    35 & 4 & 0.333 & 0.750 \\
    36 & 5 & 0.255 & 0.400 \\
    \midrule
    Overall & 17 & \textbf{0.206} & \textbf{0.353} \\
    \bottomrule
  \end{tabular}
\end{table}

The overall mIoU of 0.206 is notably higher than the 0.113 reported below
for annotated queries, but the comparison is not straightforward. The VLM
produces only 4--5 coarse queries per video---descriptions like ``crowd
dancing'' or ``people inside a train''---whereas the annotated ground truth
contains 10--19 fine-grained event descriptions per video. CLIP's joint
image-text space was trained on similarly coarse web captions, so the VLM
queries land closer to the model's natural operating regime. Video~35
illustrates both the strength and the limit of this approach: R@0.3 of
0.75 means three of the four VLM queries retrieved a reasonable window, but
the missing query corresponded to the advertisement text reveal at the end
of the video, which requires reading on-screen text rather than recognising
visual content.

\subsection*{Fine-grained evaluation: annotated events as queries}

To stress-test the retrieval with precise, human-authored descriptions, we
ran all four method variants on the full 57 merged ground-truth events.

\begin{table}[h]
  \centering
  \caption{Moment retrieval on all 57 annotated ground-truth events.}
  \label{tab:mr}
  \begin{tabular}{lrrr}
    \toprule
    Method & mIoU & R@0.3 & R@0.5 \\
    \midrule
    CLIP-B/32            & 0.089 & 0.140 & 0.070 \\
    CLIP-B/32 + fusion   & 0.082 & 0.105 & 0.053 \\
    \textbf{CLIP-L/14}   & \textbf{0.113} & \textbf{0.211} & 0.088 \\
    CLIP-L/14 + fusion   & 0.107 & 0.193 & \textbf{0.105} \\
    \bottomrule
  \end{tabular}
\end{table}

CLIP-L/14 is the clear winner, gaining roughly 27\,\% in mIoU and 50\,\%
in R@0.3 over the smaller backbone at approximately three times the
inference cost---still under 20 seconds per video on our hardware.
Paraphrase fusion, however, consistently hurts performance: both backbones
lose mIoU and R@0.3 when fusion is applied, though R@0.5 increases
slightly for CLIP-L/14 (0.105 vs.\ 0.088). We interpret this as fusion
narrowing the predicted windows---which occasionally helps at the strict
0.5 threshold---while blurring the most discriminative text embedding,
which hurts coarse alignment. A cross-attention architecture such as
CG-DETR, which can process multiple query phrasings jointly, would likely
profit from the same paraphrase inputs.

The per-category breakdown reveals where the model succeeds and where it
breaks down fundamentally. Performance events---dancing, playing
instruments---account for 29 of the 57 ground-truth events and achieve a
mean IoU of 0.174, the highest of any category, because CLIP's training
data contains abundant examples of musicians and choreography. Beekeeping
events (14 of 57) score 0.070: CLIP knows what bees and hives look like,
but cannot distinguish ``the beekeeper shows a frame'' from ``the beekeeper
opens a hive''. Speech and dialogue events (9 of 57) score exactly 0.000
across all four methods. This is not a failure of any particular design
choice---it is a fundamental limitation: CLIP has no access to audio or
text within the video frame, so descriptions that depend on what someone
is saying cannot be matched by visual similarity alone. The relationship
between saliency and IoU is non-monotonic and revealing. The most important
events (saliency $+2$, 21 events) achieve mean IoU 0.151; events rated
$+1$ (14) drop to 0.091; events rated 0 (11) drop further to 0.027;
then $-1$ events (8) recover to 0.127. The dip at zero and $+1$ captures
exactly the transitional moments our annotations describe as ``contextually
necessary but not visually distinctive''---a conductor stepping forward, a
crowd beginning to record on their phones. Those events are hardest to
retrieve because their visual signature is generic.

%% -------------------------------------------------------
\section{Video Summary Generation}
\label{sec:summary}
%% -------------------------------------------------------

For each video we selected the predicted windows from CLIP-L/14, which
achieved the best overall mIoU. We then applied a short post-processing
pass: dropping any segment shorter than 1.5\,s, merging segments separated
by a gap of 1\,s or less, sorting the result chronologically, and
concatenating the clips into a single MP4 using the OpenCV writer (with an
ffmpeg fallback where available).

\begin{table}[h]
  \centering
  \caption{Generated summary statistics.}
  \label{tab:summaries}
  \begin{tabular}{lrrr}
    \toprule
    Video & Clips & Summary & Source \\
    \midrule
    33 &  9 & 31\,s & ${\approx}$7:30 \\
    34 &  9 & 28\,s & ${\approx}$2:25 \\
    35 &  8 & 28\,s & ${\approx}$2:25 \\
    36 & 11 & 40\,s & ${\approx}$4:00 \\
    \bottomrule
  \end{tabular}
\end{table}

The summaries are concise---7 to 20\,\% of the source duration---because
we clip each segment to its predicted window without adding any padding.
Relevance is strong for visually distinctive content: videos~33--35
produce summaries dominated by dancing and instrument-playing segments,
which are exactly the moments CLIP retrieves most accurately. The
beekeeping summary is weaker on relevance because CLIP's generic picks
(hands holding objects, white protective suit) correspond to visually
similar frames throughout the video rather than the specific named
demonstrations the beekeeper is giving. Coverage suffers predictably from
the speech problem: counting only events with IoU $\geq 0.3$, the summaries
capture 21\,\% of all ground-truth events, and none of the nine
speech-only events appear in any summary. Coherence is maintained locally
by the chronological ordering and the gap-merging step, but global
coherence in video~36 breaks down because the beekeeper's explanations
build on one another cumulatively, and extracting isolated visual clips
loses the thread of his argument. The subjectivity disagreement has a
direct trace here: because the two annotators assigned opposite saliency
rankings to the same events in video~36, the merged ground truth treats
narrative speech events and visual demonstrations as equally important.
CLIP retrieves the visually distinctive ones and discards the rest,
producing a summary that satisfies neither annotator's conception of
what the video is fundamentally about.

%% -------------------------------------------------------
\section{Critical Reflections and Possible Improvements}
%% -------------------------------------------------------

The result that stands out most from this assignment is not the IoU number
itself but what drives it: the annotation process introduces more variance
than any modelling decision we made. Our two annotators agreed well on
boundaries (mIoU 0.65 on average) but substantially less on event count
and saliency. Every IoU figure reported in Section~\ref{sec:mr} is measured
against a merged ground truth that absorbs that disagreement, so the numbers
should be read with that uncertainty in mind.

The single most impactful technical improvement would be adding automatic
speech recognition. All nine speech-and-dialogue events in our ground truth
achieve zero IoU under every method we tested, not because the retrieval
architecture is wrong but because CLIP has no way to process spoken
content. A lightweight ASR model such as Whisper-tiny (80\,MB) would
allow moment-retrieval queries to be matched against a transcript, and we
expect this one change would have a larger effect on summary quality than
any visual backbone upgrade.

Two improvements that operate within the existing visual pipeline are also
worth pursuing. Paraphrase fusion did not help in our experiments, but we
used only three static rephrasings. A model that generates
semantically diverse paraphrases---rather than fixed string templates---might
provide a genuine signal boost, especially for beekeeping events where the
annotated descriptions use specialist vocabulary that CLIP's training data
may not associate strongly with the corresponding images. The saliency score
could also be reintroduced at summary-generation time, not as an annotation
metric but as a weighting signal produced by a small language model given
the event description and the retrieved clip. Selecting clips by predicted
importance rather than treating all retrieved windows equally would give the
summary a more coherent narrative arc, particularly for video~36 where the
cumulative structure of the beekeeper's explanation matters most.

Finally, the CG-DETR comparison we planned would have been instructive.
Published results on QVHighlights report mIoU around 0.40 for CG-DETR,
compared to 0.11 for our best CLIP variant. We implemented the integration code in
\texttt{pipeline/\allowbreak moment\_retrieval.py}, but the official
pre-trained checkpoint is no longer publicly available, so we
cannot report actual numbers. Once a working checkpoint is released, the
comparison would be a direct test of whether the 25--30 IoU point gap
holds on our informal, annotator-produced ground truth.

%% -------------------------------------------------------
\section{Conclusion}
%% -------------------------------------------------------

We built a complete event-detection-to-summary-video pipeline using
zero-shot pretrained models, evaluated it against our own annotations with
three inter-annotator agreement metrics and IoU across two moment-retrieval
architectures, and ran failure analysis broken down by event category and
saliency level. The dominant finding is that subjectivity in annotation
is the largest single source of variance in the whole pipeline: the two
annotators agreed on boundaries but disagreed on what counts as an event
and how important each one is, and those disagreements propagate into every
downstream metric. The second dominant finding is that zero-shot visual
retrieval fails completely on speech events, and the practical path forward
for this kind of pipeline runs through automatic speech recognition rather
than through larger or better-tuned visual encoders.

\end{document}
